% =============================================================================
%  §3–§4
% =============================================================================
\section{Springer 解消上的凝聚层}

\subsection{Springer 解消}\label{sec:springer}

用 $\cN$ 表示 $G^\vee$ 的幂零锥。我们可以考虑 Springer 解消
\[
\tN:=G^\vee\times_{B^\vee}\frn^\vee .
\]
Springer 解消带有一个矩映射 $\mu:\tN\to\cN$，以及一个投影映射 $\tN\to G^\vee/B^\vee$。
后一个映射使 $\tN$ 成为 $G^\vee/B^\vee$ 上的向量丛。
此外，它是 $\frg^\vee\times G^\vee/B^\vee$ 的子向量丛。
对每个 $\lambda\in X^\vee$，$G^\vee/B^\vee$ 上有一个线丛 $\cO_{G^\vee/B^\vee}(\lambda)$。
我们将把 $\cO_{G^\vee/B^\vee}(\lambda)$ 到 $\tN$ 的拉回记作 $\cO_{\tN}(\lambda)$。

\begin{yizhu}[术语背景：Springer 解消]
$\tN=\{(x,gB^\vee):x\in\frn^\vee\cap\mathrm{Ad}_{g^{-1}}\frb^\vee\}$ 是「把每个幂零元素连同
一个使它落在 $\frn^\vee$ 里的 Borel 子代数一起记住」的空间，故又称 Springer 簇。
投影 $\tN\to\cN$、$(x,gB^\vee)\mapsto x$ 是双有理的、对称幂零轨道上纤维是 Springer 纤维，
其最高上同调给出 Weyl 群的 Springer 表示；
nLab 条目 \texttt{Springer correspondence}
（\url{https://ncatlab.org/nlab/show/Springer+correspondence}）有概述。
本文只用它的几何：$\tN$ 是 $G^\vee/B^\vee$ 上的向量丛，且 $\cO_{\tN}(\lambda)$ 可沿 $\mu$ 推出。
\end{yizhu}

\begin{lemma}\label{lem:gen}
范畴 $D^b\Coh^{G^\vee}(\tN)$ 由下列各类对象生成：
\begin{enumerate}[label=(\arabic*),leftmargin=2.2em,itemsep=1pt]
\item 线丛 $\cO_{\tN}(\lambda)$，其中 $\lambda\in X^\vee$；
\item 形如 $V\otimes\cO_{\tN}(\lambda)$ 的对象，其中 $V\in\Rep(G^\vee)$ 且 $\lambda\in X^\vee_+$。
\end{enumerate}
\end{lemma}

我们用 $\hN$ 表示 $\tN$ 在映射 $\frg^\vee\times G^\vee/U^\vee\to G^\vee/B^\vee$ 下的原像。
我们还将考虑 $\cX:=\mathrm{Spec}\bigl(\cO(G^\vee/U^\vee)\bigr)$，即 $G^\vee/U^\vee$ 的仿射完备化。
注意 $G^\vee/U^\vee$ 是 $\cX$ 中的开集。
我们用 $\ddX$ 表示 $G^\vee/U^\vee$ 在 $\cX$ 中的补集。
不难验证，$G^\vee$ 在 $G^\vee/U^\vee$ 上作用的\emph{无穷小万有稳定化子}就是 $\hN$。
然后我们定义 $\hNX$ 为 $G^\vee$ 在 $\cX$ 上作用的无穷小万有稳定化子。
它是 $\frg^\vee\times\cX$ 的闭子概型，并且以 $\hN$ 为其开子簇。
我们用下图总结上述各对象以及它们之间的映射：

\begin{equation}
\begin{tikzcd}[column sep=4.2em, row sep=2.4em]
\tN & \hN \arrow[l,"T^\vee\text{-torsor}"'] \arrow[d,hook]
   & \hNX \arrow[l,hook',"\text{open}"] \arrow[d,hook] \\
\frg^\vee\times G^\vee/B^\vee
   & \frg^\vee\times G^\vee/U^\vee \arrow[l]
   & \frg^\vee\times\cX \arrow[l]
\end{tikzcd}
\end{equation}
图中的竖直映射都是闭子概型的包含。

\begin{yizhu}[术语背景：无穷小万有稳定化子]
设群 $G^\vee$ 作用在 $\cX$ 上。「无穷小万有稳定化子」是把每个点 $x\in\cX$ 的稳定化子
$\mathrm{Stab}(x)\subset G^\vee$ 的 Lie 代数收集起来得到的空间
$\{(x,\xi):\xi\in\mathrm{Lie}\,\mathrm{Stab}(x)\}\subset\frg^\vee\times\cX$；
它自动是闭子概型。对 $G^\vee/U^\vee\subset\cX$ 这个开集，它恰好是 $\hN$（$\tN$ 沿 $G^\vee/U^\vee\to G^\vee/B^\vee$ 的原像）；
而在完备化 $\cX$ 上，它变成 $\hNX$，比 $\hN$ 多出边界 $\ddX$ 上的部分。
图 (4) 就是在说：$\hN$ 是 $\hNX$ 的开子簇，$\hN$ 又是 $\tN$ 的 $T^\vee$-主丛。
\end{yizhu}

我们用 $\Coh^{G^\vee\times T^\vee}_{\mathrm{free}}(\hNX)$ 表示 $\Coh^{G^\vee\times T^\vee}(\hNX)$ 中
由\emph{自由}凝聚层构成的满加法子范畴，即形如 $V\otimes\cO_{\hNX}$ 的层，其中 $V\in\Rep(G^\vee\times T^\vee)$。
我们还可以考虑 $K^b\Coh^{G^\vee\times T^\vee}_{\mathrm{free}}(\hNX)$ 的满子范畴
$K^b\Coh^{G^\vee\times T^\vee}_{\mathrm{free}}(\hNX)_{\ddX}$，
它由「上同调在投影 $\hNX\to\cX$ 下支撑在 $\ddX$ 的原像上（集合论意义下）」的对象构成。

\begin{yizhu}[原文含糊：为什么用 $K^b$ 而不用 $D^b$]
$\hNX$ 不是正则的，$\Coh$ 的有界导出范畴与同伦范畴未必一致，
所以这里只在「自由层」这个加法范畴上取同伦范畴 $K^b$；
而自由层都是 $V\otimes\cO_{\hNX}$ 这种「常系数」对象，其上同调行为完全由 $V$ 决定，
用 $K^b$ 就足以做后面的消去。命题 3.1.2 说的正是：把边界 $\ddX$ 上的部分杀掉之后，
这个 $K^b$ 与 $\tN$ 上的 $D^b\Coh^{G^\vee}$ 是同一个东西。
\end{yizhu}

\begin{proposition}\label{prop:Kquot}
存在函子
\[
K^b\Coh^{G^\vee\times T^\vee}_{\mathrm{free}}(\hNX)\longrightarrow D^b\Coh^{G^\vee}(\tN),
\]
它唯一地通过三角范畴的等价
\[
K^b\Coh^{G^\vee\times T^\vee}_{\mathrm{free}}(\hNX)\big/
K^b\Coh^{G^\vee\times T^\vee}_{\mathrm{free}}(\hNX)_{\ddX}
\;\cong\; D^b\Coh^{G^\vee}(\tN)
\]
分解。

\begin{proof}
这个函子很容易构造。它由拉回函子
\[
D^b\Coh^{G^\vee\times T^\vee}(\hNX)\longrightarrow D^b\Coh^{G^\vee\times T^\vee}(\hN)
\]
连同等价\footnote{这个等价可以看成底栈的同构 $\hN/(G^\vee\times T^\vee)\cong\tN/G^\vee$。}
\[
D^b\Coh^{G^\vee\times T^\vee}(\hN)\cong D^b\Coh^{G^\vee}(\tN)
\]
诱导。验证所求函子是等价，基本上只是同调代数与代数群表示论的一道习题。
\end{proof}
\end{proposition}

\subsection{函子 I 的构造}\label{sec:F1}

我们先定义一个函子
\begin{align*}
F_1:\;&\Rep(G^\vee\times T^\vee)\longrightarrow \PWak,\\
&V\otimes\bbk_{T^\vee}(\lambda)\longmapsto Z\bigl(\Sat(V)\bigr)\conv^{I}\mathcal{W}_\lambda
\end{align*}
其中 $V\in\Rep(G^\vee)$，$\lambda\in X^\vee$。\footnote{由于我们对 Wakimoto 层的表述方式，这个函子是否真正函子性并不显然。更准确地说，可以定义一个函子 $\Rep(T^\vee)\to P_I$，它把 $\bbk_{T^\vee}(\lambda)$ 送到 $\mathcal{W}_\lambda$。}
由引理 2.5.4 与定理 2.5.3，这个函子良定义。

\begin{yizhu}[原文含糊：$F_1$ 的定义域为什么是 $G^\vee\times T^\vee$]
$V\otimes\bbk_{T^\vee}(\lambda)$ 是 $\Rep(G^\vee\times T^\vee)$ 中「$G^\vee$ 部分取 $V$、$T^\vee$ 部分取权 $\lambda$」的一维特征，
于是 $F_1$ 把它拆成两半：$Z(\Sat(V))$ 管 $G^\vee$（中心层），$\mathcal{W}_\lambda$ 管 $\lambda$（Wakimoto 层）。
函子性上的隐患见脚注 9：$\mathcal{W}_\lambda$ 只对 $\lambda$ 逐项定义，$\lambda\mapsto\mathcal{W}_\lambda$ 本身未必是函子，
但 $F_1$ 在 $\Rep(G^\vee\times T^\vee)$ 上仍然良定义。
\end{yizhu}

下一条是定理 2.4.1 (3) 与命题 2.5.2 的容易推论。

\begin{lemma}\label{lem:F1mon}
函子 $F_1$ 是幺半的。
\end{lemma}

\subsection{函子 II 的构造}\label{sec:F2}

接下来，我们希望把函子 $F_1$ 推广为一个幺半函子
\[
F_2:\Coh^{G^\vee\times T^\vee}_{\mathrm{free}}(\hNX)\longrightarrow \PWak .
\]
首先，我们把 $\PWak$ 换成一个（非满的）子范畴 $\mathcal{C}$，它是对称幺半的。\footnote{这多少是被如下观察所迫：$\Coh^{G^\vee\times T^\vee}_{\mathrm{free}}(\hNX)$ 是对称幺半的，而 $\PWak$ 不是。}
于是可以验证 $F_1$ 通过 $\mathcal{C}$ 分解，得到一个函子
\[
F_1':\Rep(G^\vee\times T^\vee)\longrightarrow\mathcal{C}.
\]
由某种类似 Tannaka 的形式主义，存在一个等价
\[
H:\mathcal{C}\longrightarrow A\text{-}\mathrm{mod}^{G^\vee\times T^\vee}_{\mathrm{free}},
\]
其中 $A$ 是某个带有 $G^\vee\times T^\vee$-作用的代数。
可以验证，经由 $H$，$F_1'$ 等同于函子
\begin{align*}
\Rep(G^\vee\times T^\vee)&\longrightarrow A\text{-}\mathrm{mod}^{G^\vee\times T^\vee}_{\mathrm{free}},\\
V&\longmapsto V\otimes A .
\end{align*}
于是，把 $F_1'$ 推广为函子 $\Coh^{G^\vee\times T^\vee}_{\mathrm{free}}(\hNX)\to\PWak$，
等价于给出一个 $G^\vee\times T^\vee$-等变代数的态射
\[
\cO(\hNX)\longrightarrow A .
\]
这件事可以通过把 $\cO(\hNX)$ 看作 $\cO(\frg^\vee)\otimes\cO(\cX)$ 的商来完成。
我们先定义等变代数态射 $\cO(\frg^\vee)\to A$ 与 $\cO(\cX)\to A$。
下面简要说明它们是如何构造的：
\begin{enumerate}[label=(\arabic*),leftmargin=2.2em,itemsep=1pt]
\item 态射 $\cO(\frg^\vee)\to A$ 实际上是附近循环的单值化自同态的「对数」。
\item 态射 $\cO(\cX)\to A$ 由定理 2.5.3 给出的中心层的 Wakimoto 滤过决定。
\end{enumerate}
于是不难验证，态射 $\cO(\frg^\vee)\otimes\cO(\cX)\to A$ 唯一地通过所求的态射 $\cO(\hNX)\to A$ 分解。
函子 $F_2$ 于是是如下函子的复合：
\begin{align*}
F_2:\;&\Coh^{G^\vee\times T^\vee}_{\mathrm{free}}(\hNX)\longrightarrow
A\text{-}\mathrm{mod}^{G^\vee\times T^\vee}_{\mathrm{free}}\cong\mathcal{C}\hookrightarrow\PWak,\\
&M\longmapsto M\otimes_{\cO(\hNX)}A .
\end{align*}

\begin{yizhu}[原文含糊：Tannaka 形式主义那一步]
这里用的是「代数表示 $\leftrightarrow$ 模范畴」的逆向 Tannaka 论证：
若一个对称幺半范畴 $\mathcal{C}$ 被一个对象 $A$ 张成（$\mathcal{C}$ 中每个对象都是 $A$ 的直和项），
则在合适的有限性条件下 $\mathcal{C}\cong A\text{-}\mathrm{mod}^{\mathrm{free}}$。
之所以要先把 $\PWak$ 换成子范畴 $\mathcal{C}$，正是因为 $\PWak$ 不是对称幺半的，
而 $\Coh_{\mathrm{free}}^{G^\vee\times T^\vee}(\hNX)$ 是——脚注 10 已经把理由说了。
\end{yizhu}

\subsection{函子 III 的构造}\label{sec:F3}

回忆命题 3.1.2：存在范畴的等价
\[
K^b\Coh^{G^\vee\times T^\vee}_{\mathrm{free}}(\hNX)\big/
K^b\Coh^{G^\vee\times T^\vee}_{\mathrm{free}}(\hNX)_{\ddX}
\;\cong\; D^b\Coh^{G^\vee}(\tN).
\]
为验证 $F_2$ 可以推广为函子
\[
D^b\Coh^{G^\vee}(\tN)\longrightarrow D^b\PWak,
\]
只需验证 $K^bF_2$ 把 $K^b\Coh^{G^\vee\times T^\vee}_{\mathrm{free}}(\hNX)_{\ddX}$
送到 $K^b\PWak$ 中的零调复形。
这可以相当显式地验证，故从略。

结论是存在函子
\begin{align*}
F:\;&D^b\Coh^{G^\vee}(\tN)\longrightarrow D^b\PWak\hookrightarrow D^bP_I,\\
&V\otimes\cO_{\tN}\longmapsto Z\bigl(\Sat(V)\bigr),\\
&\cO_{\tN}(\lambda)\longmapsto\mathcal{W}_\lambda .
\end{align*}
我们还会提到这个函子 $F$ 的一个变体，即把它与实现函子复合：
\[
\bar F:D^b\Coh^{G^\vee}(\tN)\longrightarrow D^bP_I\longrightarrow D^b(I\backslash\Fl_G).
\]
我们要强调，$\Fl_G$ 上 $I$-等变反常层的实现函子并不是等价。

\begin{keytakeaways}
本节把「$\tN$ 上的凝聚层」翻译成「$\Fl_G$ 上的层」，分三步：
$F_1$ 在表示范畴上定义（$V\otimes\bbk(\lambda)\mapsto Z(\Sat V)\conv\mathcal{W}_\lambda$），
$F_2$ 借 Tannaka 形式推广到自由凝聚层，$F_3$ 再用命题 3.1.2 的消去把边界 $\ddX$ 除掉。
三步的每一步都只用到上一节的两件工具，这是本文最漂亮的地方：
复杂的凝聚层一侧被完全压进「中心层 $+$ Wakimoto 层」这两个可计算的部件里。
\end{keytakeaways}

% =============================================================================
\section{反球范畴与 Iwahori–Whittaker 层}

我们将构造反球模的另一种范畴化。
具体地说，我们将定义 \emph{Iwahori–Whittaker 反常层}的范畴，记作 $\PIW$，
并给出一个正合、全忠实的函子
\[
\Pasph\longrightarrow \PIW .
\]
我们将证明，中心层在函子
\[
\Rep(G^\vee)\longrightarrow P_I\longrightarrow\Pasph\longrightarrow\PIW
\]
下的像就是\emph{倾斜对象}（tilting object）。

\begin{remark}
为了说清 Iwahori–Whittaker 层，我们需要离开 Betti 设定。\footnote{这并非绝对必要。我们可以在 Betti 设定下改用 Krillov 模型（\cite{Gai2}），但这类考量超出本注记的范围。}
取而代之，我们将固定 $\bbF$ 为 $\bbF_p$ 的代数闭包。
于是 $G$ 将是 $\bbF$ 上的连通约化群，同样地，所有 ind-概型都将定义在 $\bbF$ 上。
我们还将固定 $\ell\neq p$，并且我们所有的可构造层都应解释为平展 $\bbQl$-层。
至此讨论过的一切对平展层都有显然的类比。
\end{remark}

\begin{yizhu}[原文疑误：Krillov]
脚注 11 中的「Krillov model」应为「Kirillov model」，即 Whittaker 模型在 $GL_2$ 情形下的经典名称。
译文照录原文，此处标出。
\end{yizhu}

\begin{yizhu}[术语背景：为何要离开 Betti 设定]
Artin–Schreier 层是特征 $p$ 的现象：它来自 $x\mapsto x^p-x$ 的伽罗瓦覆叠，
在复数（Betti）设定下没有对应物，只能改用指数层的类比或 Kirillov 模型。
因此第 4 节把工作域换成 $\overline{\bbF}_p$，并把所有层换成平展 $\bbQl$-层；
这也解释了为什么第 2、3 节一直在 $\bbC$ 上工作，而第 4 节一开始就要「换场景」。
\end{yizhu}

\subsection{热身：Whittaker 层}\label{sec:whittaker}

首先，我们解释 Whittaker 层的一个简化版本。
设 $X$ 是一个 $H$-簇，其中 $H$ 是代数群（不必约化）。
用 $a:H\times X\to X$ 表示作用映射。固定一个同态 $\chi:H\to\Gm$。
人们可以考虑 $\Ga$ 上的 Artin–Schreier 局部系统 $\mathcal{L}_{AS}$。\footnote{具体地说，选取一个本原 $p$ 次单位根 $\zeta\in\bbQl$。在由 $x\mapsto x^p-x$ 定义的伽罗瓦覆叠 $\Ga\to\Ga$ 下，常层的推前有一个直和项，$\bbZ/p\bbZ$ 在其上通过特征 $[n]\mapsto\zeta^n$ 作用。}
拉回 $\chi^*\mathcal{L}_{AS}$ 是 $H$ 上的乘法局部系统。
于是人们可以考虑 $X$ 上可构造层的 $(H,\chi^*\mathcal{L}_{AS})$-等变导出范畴，
记作 $D^b_c(H\backslash_{\chi^*\mathcal{L}_{AS}}X)$。

直观上（但并不精确\footnote{要使这一点精确，必须把由乘法映射 $H^k\times X\to X$ 给出的高阶相容关系也纳入进来，参见 \cite{Gai2}。}），
$D^b_c(H\backslash_{\chi^*\mathcal{L}_{AS}}X)$ 的对象由层 $\cF\in D^b_c(X)$
连同自然同构
\[
a^*\cF\cong\chi^*\mathcal{L}_{AS}\boxtimes\cF
\]
组成。

\begin{yizhu}[术语背景：Artin–Schreier 局部系统]
$\mathcal{L}_{AS}$ 是 $\Ga$ 上秩一、$p$-挠的平展局部系统，
其单值化由 $x\mapsto x^p-x$ 的覆叠给出；它是 $\Ga$ 上唯一（在差一个平移意义下）
满足「乘性」的局部系统，因而可以沿 $\chi$ 拉回到 $H$ 上。
Whittaker 层就是「沿着 $\chi$ 扭曲等变」的层：它把等变条件里的常层换成了 $\chi^*\mathcal{L}_{AS}$。
关键性质是 $H^\bullet(\Ga,\mathcal{L}_{AS})=0$（见下文），
这使 Whittaker 侧的上同调大量消没，从而可以像「局部化」一样把层论问题压到很小的支撑上。
\end{yizhu}

Whittaker 层往往满足强的上同调消没性质，其动机是如下基本性质：
\[
H^{\bullet}(\Ga,\mathcal{L}_{AS})\cong H^{\bullet}_c(\Ga,\mathcal{L}_{AS})=0 .
\]
因此，Whittaker 层往往受到很紧的控制：

\begin{example}
设 $\Ga$ 通过 $t\cdot[x:y]=[x+yt:y]$ 作用在 $\bbP^1$ 上。
那么 $\bbP^1$ 上有两条 $\Ga$-轨道：一条同构于 $\bbA^1$，另一条同构于一个点。
点这条轨道太小了，无法承载 Whittaker 等变层，因为点上的任何层都是 $\Ga$-等变的。
于是
\[
D^b_c(\Ga\backslash_{\chi}\bbP^1)\cong D^b_c(\Ga\backslash_{\chi}\bbA^1).
\]
然后不难验证，$D^b_c(\Ga\backslash_{\chi}\bbA^1)\cong D^b(\bbQl\text{-}\mathrm{mod})$，
它由在 $0\in\bbA^1$ 处取茎所诱导。
\end{example}

\begin{yizhu}[原文含糊：例 4.1.1 的「取茎」]
$0\in\bbA^1$ 是 $\Ga$-作用的不动点，也是 $\bbA^1$ 上唯一使 $\chi$-等变条件非平凡的点；
在 $0$ 处取茎的函子 $D^b_c(\Ga\backslash_\chi\bbA^1)\to D^b(\bbQl\text{-}\mathrm{mod})$
是全忠实且本质满的（等价），因为它把整个范畴「压」到那条不动轨道上。
这也是 Whittaker 层「局部化」的几何雏形：所有信息都集中到 $\chi$ 非平凡的轨道上。
\end{yizhu}

\subsection{Iwahori–Whittaker 范畴}\label{sec:iwcat}

设 $I^+_u$ 是 $I^+$ 的 pro-幺幂根，即 $U^+$ 在自然映射 $LG\to G$ 下的原像。
我们将固定 $U^+$ 的定位（pinning），以得到同构 $U^+/[U^+,U^+]\cong\prod\Ga$。
定义一个群同态，它是如下复合：
\[
\chi:\;I^+_u\longrightarrow U^+\longrightarrow U^+/[U^+,U^+]\cong\prod\Ga\xrightarrow{\ \mathrm{sum}\ }\Ga .
\]

我们定义 \emph{Iwahori–Whittaker 层}的范畴，记作 $D^b_{\IW}(\Fl_G)$，
为 $\Fl_G$ 上 $\bbQl$-层的 $(I^+_u,\chi^*(\mathcal{L}_{AS}))$-等变导出范畴。
$D^b_c(\Fl_G)$ 上的反常 $t$-结构限制为 $D^b_{\IW}(\Fl_G)$ 上的 $t$-结构。
我们用 $\PIW$ 表示反常 Iwahori–Whittaker 层构成的范畴。
$D^b_c(I\backslash\Fl_G)$ 在 $D^b_{\IW}(\Fl_G)$ 上有一个右作用，
其定义与 $D^b_c(I\backslash\Fl_G)$ 上的卷积类似。即存在双函子
\[
(-)\conv^{I}(-):\;D^b_{\IW}(\Fl_G)\times D^b_c(I\backslash\Fl_G)\longrightarrow D^b_{\IW}(\Fl_G).
\]

\begin{yizhu}[术语背景：pro-幺幂根与 Whittaker 特征]
$I^+$ 是 $L^+G$ 中对应于 $B^+$ 的 Iwahori 子群，$U^+\subset G$ 是其幺幂部分。
$I^+_u$ 是 $I^+$ 的 pro-幺幂根，可以看作 $U^+$ 在环路群里的「$\mathcal{O}$-点」版本。
由于 $U^+$ 幺幂，$U^+/[U^+,U^+]\cong\prod\Ga$，而 $\mathrm{sum}$ 把它映到 $\Ga$，
再用 Artin–Schreier 局部系统扭出特征 $\chi$：这是一维 Whittaker 特征在环路群上的提升。
「Iwahori–Whittaker」这个名字正是「Iwahori 等变」与「Whittaker 特征」两件事的叠加。
\end{yizhu}

\subsection{最高权结构}\label{sec:hw}

$\Fl_G$ 上的 $I^+_u$-轨道由 $\Wex$ 参数化。存在双射
\begin{align*}
X^\vee&\longleftrightarrow \Wf,\\
\lambda&\longmapsto w_\lambda .
\end{align*}
用 $\Fl^{\IW}_{G,\lambda}$ 表示对应于 $w_\lambda$ 的 $I^+_u$-轨道，
并用 $j^{\IW}_\lambda:\Fl^{\IW}_{G,\lambda}\hookrightarrow\Fl_G$ 表示嵌入。

$\Fl_G$ 上只有那些由 $w\in\Wf$ 指标化的 $I^+_u$-轨道才承载非零的 Whittaker 局部系统。
事实上，对每个 $\lambda\in X^\vee$，$\Fl^{\IW}_{G,\lambda}$ 上有唯一的秩一 Whittaker 局部系统，
记作 $\mathcal{L}_{\chi,\lambda}$。对每个 $\lambda\in X^\vee$，我们定义层
\[
\Delta^{\IW}_\lambda=j^{\IW}_{\lambda!}\mathcal{L}_{\chi,\lambda}\bigl[\dim\Fl^{\IW}_{G,\lambda}\bigr],
\qquad
\nabla_\lambda=j^{\IW}_{\lambda*}\mathcal{L}_{\chi,\lambda}\bigl[\dim\Fl^{\IW}_{G,\lambda}\bigr].
\]
这些层分别称为\emph{标准层}与\emph{余标准层}。
它们自然地是 Iwahori–Whittaker 反常层。
典范态射 $\Delta^{\IW}_\lambda\to\nabla^{\IW}_\lambda$ 的像是单 Iwahori–Whittaker 反常层，
记作 $\mathrm{IC}^{\IW}_\lambda$。

\begin{proposition}\label{prop:hw}
范畴 $\PIW$ 自然地具有最高权范畴的结构，其权偏序集为 $X^\vee$。
\end{proposition}

\begin{yizhu}[术语背景：最高权范畴]
最高权范畴（highest weight category）是指一个 Abel 范畴，其单对象按偏序集 $\Lambda$ 指标化，
每个单对象 $L_\lambda$ 有一个射影覆盖 $P_\lambda$ 使得 $\mathrm{Hom}(P_\lambda,L_\mu)\neq0$ 仅当 $\mu\le\lambda$，
并且每个 $P_\lambda$ 都有「标准层 $\Delta_\lambda$」作为滤过商——即所谓「$\Delta$-滤过」。
它给出标准层、余标准层、倾斜对象这一整套机器（Cline–Parshall–Scott）。
「权偏序集为 $X^\vee$」是说这里的「权」就是支配权 $\lambda$，与 $\Fl_G$ 的 $I^+_u$-轨道一一对应。
\end{yizhu}

命题 4.3.1 的关键观察是：$\PIW$ 容许倾斜层。
具体地说，对每个 $\lambda\in X^\vee$，存在唯一的不可分解 Iwahori–Whittaker 倾斜反常层
$\mathcal{T}^{\IW}_\lambda$。

\begin{corollary}\label{cor:realiz}
实现函子诱导范畴的等价
\[
D^b\PIW\;\cong\;D^b_{\IW}(\Fl_G).
\]
\end{corollary}

\subsection{Iwahori–Whittaker 平均}\label{sec:av}

定义 \emph{Iwahori–Whittaker 平均函子}
\begin{align*}
\AvIW:\;&D^b(I\backslash\Fl_G)\longrightarrow D^b_{\IW}(\Fl_G)\\
&\cF\longmapsto\Delta^{\IW}_0\conv^{I}\cF .
\end{align*}

\begin{theorem}[Arkhipov–Bezrukavnikov, \cite{ArBe}]\label{thm:av}
\begin{enumerate}[label=(\arabic*),leftmargin=2.2em,itemsep=1pt]
\item 函子 $\AvIW$ 是反常 $t$-正合的。
\item $\AvIW$ 到反常层上的限制 $P_I\to\PIW$ 通过一个全忠实函子 $\Pasph\to\PIW$ 分解。
\end{enumerate}
\end{theorem}

\begin{proof}[定理 4.4.1 的证明梗概]
由一般的几何考量并考察轨道，可以证明
\begin{equation}
\AvIW(\mathrm{IC}_w)=0
\end{equation}
除非 $w\in\Wf$。此外，对任意 $w\in\Wex$，存在同构
\begin{equation}
\AvIW(\Delta_w)\cong\Delta^{\IW}_\lambda,
\qquad
\AvIW(\nabla_w)\cong\nabla^{\IW}_\lambda,
\end{equation}
其中 $\lambda\in X^\vee$ 是使得 $W\cdot w=W\cdot w_\lambda$ 的唯一元素。

由基本的同调代数，可以从 (5) 与 (6) 推出 $\AvIW$ 是反常 $t$-正合的。
此外，由 (5)，$\AvIW$ 将通过反球 Hecke 范畴分解。

为证明诱导出的函子 $\AvIW:\Pasph\to\PIW$ 全忠实，人们构造 $\AvIW$ 的一个截面。
直观上，这通过「从 $I\cap I^+_u$-等变平均到 $I$-等变，取反常上同调，
再把商函子作用到反球 Hecke 范畴」来完成。
\end{proof}

\begin{yizhu}[原文含糊：式 (5) 的几何直觉]
$\mathrm{IC}_w$ 是层片 $\Fl_{G,w}$ 的相交上同调复形。
$\AvIW$ 先把它与 $\Delta^{\IW}_0$（支撑在单位轨道上、带 Whittaker 局部系统）卷积，
再取 $I^+_u$-不变部分；若 $w\notin\Wf$，则 $\Fl_{G,w}$ 的闭包在某个正根方向上是「可缩」的，
Whittaker 局部系统沿该方向的上同调消没（即 $H^\bullet(\Ga,\mathcal{L}_{AS})=0$ 的几何版本），于是整项消失。
反之若 $w=w_\lambda\in\Wf$，卷积只剩下该层片上的标准层，且 Whittaker 局部系统恰好给出 $\mathcal{L}_{\chi,\lambda}$。
式 (6) 就是说：$\AvIW$ 把标准层送到标准层、余标准层送到余标准层，且下标由 $W$-轨道决定。
\end{yizhu}

\subsection{中心层与 Iwahori–Whittaker 倾斜层}\label{sec:tilt}

引入 Iwahori–Whittaker 层的主要结果与动机来自下述定理。

\begin{theorem}[Arkhipov–Bezrukavnikov, \cite{ArBe}]\label{thm:tilt}
考虑函子的复合
\[
Z^{\IW}:\;\Rep(G^\vee)\xrightarrow{\ Z\ }P_I\xrightarrow{\ \AvIW\ }\PIW .
\]
对一切 $V\in\Rep(G^\vee)$，反常层 $Z^{\IW}(V)$ 是倾斜的。
此外，对任意 $\lambda\in X^\vee$，
\[
\bigl(Z^{\IW}(V):\Delta^{\IW}_\lambda\bigr)
=\bigl(Z^{\IW}(V):\nabla^{\IW}_\lambda\bigr)
=\dim V_\lambda .
\]
\end{theorem}

\begin{yizhu}[原文疑误]
原文此式左端写作 $(Z^{\IW}(W):\Delta^{\IW}_\lambda)$，而右端、以及同一句中的量词都写作 $V$；
此处的 $W$ 显然是 $V$ 之误（$W$ 在全文别处只指 Weyl 群）。译文按 $V$ 译出，此处标出。
\end{yizhu}

\begin{yizhu}[术语背景：倾斜层与重数]
「倾斜」指对象同时具有标准滤过与余标准滤过（这里对 $\PIW$ 而言）。
$(M:\Delta_\lambda)$ 表示 $M$ 的标准滤过中 $\Delta_\lambda$ 出现的重数；
对倾斜对象，这个重数也等于余标准滤过中的重数，且等于 $M$ 在 Grothendieck 群中的坐标。
定理 4.5.1 因此是说：$Z^{\IW}$ 把不可约表示 $L_\lambda$ 送到「重数恰好为 $\dim V_\lambda$」的倾斜层，
即几何侧完美地实现了表示论里的权重数公式。
\end{yizhu}

\begin{proof}[定理 4.5.1 的证明梗概]
证明的基本想法是不断归约，直到剩下的情形可以显式手工验证。步骤如下：
\begin{enumerate}[label=(\arabic*),leftmargin=2.2em,itemsep=1pt]
\item 若 $V,V'\in\Rep(G^\vee)$ 使得 $Z^{\IW}(V)$ 与 $Z^{\IW}(V')$ 都是倾斜的，
      则 $Z^{\IW}(V\otimes V')$ 也是倾斜的。
      这基本上由 $D^b(I\backslash\Fl_G)$ 中卷积的性质推出。
\item 若定理对半单群 $G$ 成立，则它对约化群 $G$ 也成立。
\item 由 (1) 与 (2)，只需假设 $G$ 半单、且 $V$ 是最高权为极小（miniscule）或拟极小（quasi-miniscule）的单 $G^\vee$-模。
      极小的情形是相当直接的算；而拟极小的情形要困难得多，需要研究 $P_I$ 的某些商。
\end{enumerate}
\end{proof}

\begin{keytakeaways}
本节把「反球 Hecke 范畴」与「Iwahori–Whittaker 范畴」接上：
平均函子 $\AvIW$ 杀掉所有非反球的 $\mathrm{IC}_w$（式 (5)），把标准层送到标准层（式 (6)），
从而给出全忠实函子 $\Pasph\hookrightarrow\PIW$（定理 4.4.1）。
再复合上中心函子 $Z$，就得到定理 4.5.1：中心层在 $\PIW$ 里是倾斜层，且重数就是权重数。
这一步是把「$\Fl_G$ 上的层」变成「$\PIW$ 里的倾斜层」的关键跳板。
\end{keytakeaways}
